\uuid{AsfT}
\exo7id{1574}
\titre{exo7 1574}
\auteur{legall}
\organisation{exo7}
\datecreate{1998-09-01}
\isIndication{false}
\isCorrection{false}
\chapitre{Réduction d'endomorphisme, polynôme annulateur}
\sousChapitre{Polynôme annulateur}
\cpge{mpsi-mp2i.calcul-matriciel-systemes.e, pcsi.determinants.c, ptsi.determinants.a, mpsi-mp2i.polynomes-fractions-rationnelles.c, pcsi.polynomes.c, ptsi.polynomes.c, bcpst1.matrices, mp-mpi.reduction-endomorphismes.g, pc.complements-espaces-vectoriels.d, psi.complements-espaces-vectoriels.d, mp-mpi.reduction-endomorphismes.c, pc.reduction.b, psi.reduction-endomorphismes.b, pt.reduction-endomorphismes.b}
\module{Algèbre}
\niveau{L2}
\difficulte{}

\contenu{
\question{
Soit $  A\in M_3 (\R )  $ la matrice $  \begin{pmatrix} 0 & 1 & 1 \cr 
                                      1 & 0 & 1 \cr
                         0 & 0 & 1\cr \end{pmatrix} .$ 
 Calculer le polyn\^ome minimal de $  A  .$ En d\'eduire $  A^{-1}  ,$ $  A^{3}  $ et $  
 A^{5}  .$
}
}
